%   `Ueber die Anzahl der Primzahlen unter einer gegebenen Gr"osse.'
%
%   By Bernhard Riemann (1826--1866)
%
% Transcribed in LaTeX by
% 
%    David R. Wilkins
%    School of Mathematics, Trinity College, Dublin 2, Ireland
%    dwilkins@maths.tcd.ie
%
% Source:
%
%    `Bernhard Riemann's Gesammelte Mathematische Werke und
%    Wissenschaftlicher Nachlass', herausgegeben under Mitwirkung
%    von Richard Dedekind, von Heinrich Weber.
%    (Leipzig: B. G. Teubner 1892).  145--153.
% 
% David R. Wilkins
% 16th June, 1998.

\documentclass[a4paper,12pt]{article}
\usepackage[T1]{fontenc}
\usepackage[german]{babel}

\begin{document}
\thispagestyle{empty}
\begin{center}
\Large\bfseries
Ueber die Anzahl der Primzahlen unter einer gegebenen Gr\"{o}sse.\\[12 pt]
Bernhard Riemann\\[12 pt]
[Monatsberichte der Berliner Akademie, November 1859.]\\[24 pt]
\large\mdseries
Transcribed by D. R. Wilkins\\[12 pt]
Preliminary Version: December 1998
\end{center}

\newpage
\setcounter{page}{1}


\title{Ueber die Anzahl der Primzahlen unter einer gegebenen Gr\"{o}sse.}
\author{Bernhard Riemann}
\date{[Monatsberichte der Berliner Akademie, November 1859.]}

\maketitle

Meinen Dank f\"{u}r die Auszeichnung, welche mir die Akademie
durch die Aufnahme unter ihre Correspondenten hat zu Theil werden
lassen, glaube ich am besten dadurch zu erkennen zu geben, dass ich
von der hierdurch erhaltenen Erlaubniss baldigst Gebrauch mache durch
Mittheilung einer Untersuchung \"{u}ber die H\"{a}ufigkeit der
Primzahlen; ein Gegenstand, welcher durch die Interesse, welches
\emph{Gauss} und \emph{Dirichlet} demselben l\"{a}ngere Zeit
geschenkt haben, einer solchen Mittheilung vielleicht nicht ganz
unwerth erscheint.

Bei dieser Untersuchung diente mir als Ausgangspunkt die von
\emph{Euler} gemachte Bemerkung, dass das Product
\[ \prod \frac{1}{\displaystyle 1 - \frac{1}{p^s}}
   = {\textstyle\sum} \frac{1}{n^s},\]
wenn f\"{u}r $p$ alle Primzahlen, f\"{u}r $n$ alle ganzen Zahlen
gesetzt werden.  Die Function der complexen Ver\"{a}nderlichen~$s$,
welche durch diese beiden Ausdr\"{u}cke, so lange sie convergiren,
dargestellt wird, bezeichne ich durch $\zeta(s)$.  Beide convergiren
nur, so lange der reelle Theil von $s$ gr\"{o}sser als $1$ ist; es
l\"{a}sst sich indess leicht ein immer g\"{u}ltig bleibender Ausdruck
der Function finden.  Durch Anwendung der Gleichung
\[ \int\limits_0^\infty e^{-nx} x^{s-1} \,dx = \frac{\Pi(s - 1)}{n^s} \]
erh\"{a}lt man zun\"{a}chst
\[ \Pi(s - 1) \zeta(s) = \int\limits_0^\infty \frac{x^{s-1}\,dx}{e^x - 1}.\]
Betrachtet man nun das Integral
\[ \int\limits \frac{(-x)^{s-1} \,dx}{e^x - 1} \]
von $+\infty$ bis $+\infty$ positiv um ein Gr\"{o}ssengebiet erstreckt,
welches den Werth~$0$, aber keinen andern Unstetigkeitswerth der
Function unter dem Integralzeichen im Innern enth\"{a}lt, so ergiebt
sich dieses leicht gleich
\[ (e^{-\pi si} - e^{\pi si})
      \int\limits_0^\infty \frac{x^{s-1}\,dx}{e^x - 1},\]
vorausgesetzt, dass in der vieldeutigen Function
$(-x)^{s-1} = e^{(s-1)\log(-x)}$
der Logarithmus von $-x$ so bestimmt worden ist, dass er f\"{u}r ein
negatives $x$ reell wird.  Man hat daher
\[ 2 \sin \pi s \, \Pi(s - 1) \zeta(s)
   = i \int\limits_\infty^\infty \frac{(-x)^{s-1} \,dx}{e^x - 1},\]
das Integral in der eben angegebenen Bedeutung verstanden.

Diese Gleichung giebt nun den Werth der Function $\zeta(s)$ f\"{u}r
jedes beliebige complexe $s$ und zeigt, dass sie einwerthig und
f\"{u}r alle endlichen Werthe von $s$, ausser~$1$, endlich ist, so wie
auch, dass sie verschwindet, wenn $s$ gleich einer negativen geraden
Zahl ist.

Wenn der reelle Theil von $s$ negativ ist, kann das Integral, statt
positiv um das angegebene Gr\"{o}ssengebiet auch negativ um das
Gr\"{o}ssengebiet, welches s\"{a}mmtliche \"{u}brigen complexen
Gr\"{o}ssen enth\"{a}lt, erstreckt werden, da das Integral durch
Werthe mit unendlich grossem Modul dann unendlich klein ist.  Im Innern
dieses Gr\"{o}ssengebiets aber wird die Function unter dem
Integralzeichen nur unstetig, wenn $x$ gleich einem ganzen Vielfachen
von $\pm 2\pi i$ wird und das Integral ist daher gleich der Summe der
Integrale negativ um diese Werthe genommen.  Das Integral um den Werth
$n \, 2\pi i$ aber ist $= (-n \, 2\pi i)^{s-1}(-2\pi i)$, man
erh\"{a}lt daher
\[ 2 \sin \pi s \, \Pi(s - 1) \zeta(s)
   = (2\pi)^s {\textstyle\sum} n^{s-1} ((-i)^{s-1} + i^{s-1}),\]
also eine Relation zwischen $\zeta(s)$ und $\zeta(1 - s)$, welche sich
mit Benutzung bekannter Eigenschaften der Function $\Pi$ auch so
ausdr\"{u}cken l\"{a}sst:
\[ \Pi \left( \frac{s}{2} - 1 \right) \pi^{-\frac{s}{2}} \zeta(s) \]
bleibt unge\"{a}ndert, wenn $s$ in $1 - s$ verwandelt wird.

Diese Eigenschaft der Function veranlasste mich statt $\Pi(s - 1)$ das
Integral $\displaystyle \Pi \left( \frac{s}{2} - 1 \right)$ in dem
allgemeinen Gliede der Reihe $\displaystyle \sum \frac{1}{n^s}$
einzuf\"{u}hren, wodurch man einen sehr bequemen Ausdruck der Function
$\zeta(s)$ erh\"{a}lt.  In der That hat man
\[ \frac{1}{n^s} \Pi \left( \frac{s}{2} - 1 \right) \pi^{-\frac{s}{2}}
   = \int\limits_0^\infty e^{-nn\pi x} x^{\frac{s}{2} - 1} \,dx,\]
also, wenn man
\[ \sum_1^\infty e^{-nn\pi x} = \psi(x) \]
setzt,
\[ \Pi \left( \frac{s}{2} - 1 \right) \pi^{-\frac{s}{2}} \zeta(s)
   = \int\limits_0^\infty \psi(x) x^{\frac{s}{2} - 1} \,dx,\]
oder da
\[ 2 \psi(x) + 1 = x^{-\frac{1}{2}} \left(
         2\psi \left( \frac{1}{x} \right) + 1 \right),
   \mbox{ (Jacobi, Fund. S. 184)}\]
\begin{eqnarray*}
\Pi \left( \frac{s}{2} - 1 \right) \pi^{-\frac{s}{2}} \zeta(s)
   &=& \int\limits_1^\infty \psi(x) x^{\frac{s}{2} - 1} \,dx
      + \int\limits_1^\infty \psi \left( \frac{1}{x} \right)
            x^{\frac{s-3}{2}} \,dx \\
   & &+ {\textstyle\frac{1}{2}} \int\limits_0^1 \left(
            x^{\frac{s-3}{2}} - x^{\frac{s}{2} - 1}
            \right) \,dx \\
   &=& \frac{1}{s(s-1)} + \int\limits_1^\infty \psi(x) \left(
            x^{\frac{s}{2} - 1} + x^{-\frac{1+s}{2}}
            \right) \,dx.
\end{eqnarray*}

Ich setze nun $s = \frac{1}{2} + ti$ und
\[ \Pi \left( \frac{s}{2} \right) (s - 1) \pi^{-\frac{s}{2}} \zeta(s)
   = \xi(t),\]
so dass
\[ \xi(t) = {\textstyle\frac{1}{2}}
      - (tt + {\textstyle\frac{1}{4}}) \int\limits_1^\infty
         \psi(x) x^{-\frac{3}{4}}
         \cos ({\textstyle\frac{1}{2}}t \log x) \,dx \]
oder auch
\[ \xi(t)
   = 4 \int\limits_1^\infty \frac{d(x^{\frac{3}{2}} \psi'(x))}{dx}
          x^{-\frac{1}{4}} \cos ({\textstyle\frac{1}{2}}t \log x) \,dx.\]

Diese Function ist f\"{u}r alle endlichen Werthe von $t$ endlich, und
l\"{a}sst sich nach Potenzen von $tt$ in eine sehr schnell
convergirende Reihe entwickeln.  Da f\"{u}r einen Werth von $s$,
dessen reeller Bestandtheil gr\"{o}sser als $1$ ist,
$\log \zeta(s) = - \sum \log (1 - p^{-s})$
endlich bleibt, und von den Logarithmen der \"{u}brigen Factoren von
$\xi(t)$ dasselbe gilt, so kann die Function $\xi(t)$ nur
verschwinden, wenn der imagin\"{a}re Theil von $t$ zwischen
$\frac{1}{2} i$ und $-\frac{1}{2} i$ liegt.  Die Anzahl der Wurzeln
von $\xi(t) = 0$, deren reeller Theil zwischen $0$ und $T$ liegt, ist
etwa
\[ = \frac{T}{2\pi} \log \frac{T}{2\pi} - \frac{T}{2\pi};\]
denn das Integral $\int\limits d\log \xi(t)$ positiv um den Inbegriff der
Werthe von $t$ erstreckt, deren imagin\"{a}rer Theil zwischen
$\frac{1}{2} i $ und $-\frac{1}{2}i$ und deren reeller Theil zwischen
$0$ und $T$ liegt, ist (bis auf einen Bruchtheil von der Ordnung der
Gr\"{o}sse
$\displaystyle \frac{1}{T}$) gleich
$\displaystyle \left( T \log \frac{T}{2\pi} - T \right) i$;
dieses Integral aber ist gleich der Anzahl der in diesem Gebiet
liegenden Wurzeln von $\xi(t) = 0$, multiplicirt mit $2\pi i$.  Man
findet nun in der That etwa so viel reelle Wurzeln innerhalb dieser
Grenzen, und es ist sehr wahrscheinlich, dass alle Wurzeln reell sind.
Hiervon w\"{a}re allerdings ein strenger Beweis zu w\"{u}nschen; ich
habe indess die Aufsuchung desselben nach einigen fl\"{u}chtigen
vergeblichen Versuchen vorl\"{a}ufig bei Seite gelassen, da er
f\"{u}r den n\"{a}chsten Zweck meiner Untersuchung entbehrlich schien.

Bezeichnet man durch $\alpha$ jede Wurzel der Gleichung
$\xi(\alpha) = 0$, so kann man $\log \xi(t)$ durch
\[ {\textstyle\sum} \log \left( 1 - \frac{tt}{\alpha \alpha} \right)
      + \log \xi(0) \]
ausdr\"{u}cken; denn da die Dichtigkeit der Wurzeln von der
Gr\"{o}sse~$t$ mit $t$ nur wie
$\displaystyle \log \frac{t}{2\pi}$
w\"{a}chst, so convergirt dieser Ausdruck und wird f\"{u}r ein
unendliches $t$ nur unendlich wie $t \log t$; er unterscheidet sich
also von $\log \xi(t)$ um eine Function von $tt$, die f\"{u}r ein
endliches $t$ stetig und endlich bleibt und mit $tt$ dividirt f\"{u}r
ein unendliches $t$ unendlich klein wird.  Dieser Unterschied ist
folglich eine Constante, deren Werth durch Einsetzung von $t = 0$
bestimmt werden kann.

Mit diesen H\"{u}lfsmitteln l\"{a}sst sich nun die Anzahl der
Primzahlen, die kleiner als $x$ sind, bestimmen.

Es sei $F(x)$, wenn $x$ nicht gerade einer Primzahl gleich ist,
gleich dieser Anzahl, wenn aber $x$ eine Primzahl ist, um
$\frac{1}{2}$ gr\"{o}sser, so dass f\"{u}r ein $x$, bei welchem
$F(x)$ sich sprungweise \"{a}ndert,
\[ F(x) = \frac{F(x + 0) + F(x - 0)}{2}.\]

Ersetzt man nun in
\[ \log \zeta(s) = - {\textstyle\sum} \log (1 - p^{-s})
   = {\textstyle\sum} p^{-s}
     + {\textstyle\frac{1}{2} \sum} p^{-2s}
     + {\textstyle\frac{1}{3} \sum} p^{-3s} + \cdots \]
\[ p^{-s}  \mbox{ durch } s\int\limits_p^\infty x^{-s-1}\,ds,\quad
   p^{-2s} \mbox{ durch } s\int\limits_{p^2}^\infty x^{-s-1}\,ds,\ldots,\]
so erh\"{a}lt man
\[ \frac{\log \zeta(s)}{s}
   = \int\limits_1^\infty f(x) x^{-s-1} \,dx,\]
wenn man
\[ F(x) + {\textstyle\frac{1}{2}} F(x^{\frac{1}{2}})
        + {\textstyle\frac{1}{3}} F(x^{\frac{1}{3}}) + \cdots \]
durch $f(x)$ bezeichnet.

Diese Gleichung ist g\"{u}ltig f\"{u}r jeden complexen Werth $a + bi$
von $s$, wenn $a > 1$.  Wenn aber in diesem Umfange die Gleichung
\[ g(s) = \int\limits_0^\infty h(x) x^{-s} \,d \log x \]
gilt, so kann man mit H\"{u}lfe des \emph{Fourier}'schen Satzes die
Function~$h$ durch die Function~$g$ ausdr\"{u}cken.  Die Gleichung
zerf\"{a}llt, wenn $h(x)$ reell ist und
\[ g(a + bi) = g_1(b) + i g_2(b),\]
in den beiden folgenden:
\[ g_1(b)   =    \int\limits_0^\infty h(x) x^{-a} \cos (b \log x) \,d \log x,\]
\[ i g_2(b) = -i \int\limits_0^\infty h(x) x^{-a} \sin (b \log x) \,d \log x.\]

Wenn man beide Gleichungen mit
\[ (\cos (b \log y) + i \sin (b \log y)) \,db \]
multiplicirt und von $-\infty$ bis $+\infty$ integrirt, so erh\"{a}lt
man in beiden auf der rechten Seite nach dem \emph{Fourier}'schen
Satze $\pi h(y) y^{-\alpha}$, also, wenn man beide Gleichungen addirt
und mit $i y^\alpha$ multiplicirt,
\[ 2\pi i h(y) = \int\limits_{a - \infty i}^{a + \infty i} g(s) y^s \,ds,\]
worin die Integration so auszuf\"{u}hren ist, dass der reelle Theil
von $s$ constant bleibt.

Das Integral stellt f\"{u}r einen Werth von $y$, bei welchem eine
sprungweise Aenderung der Function $h(y)$ stattfindet, den Mittelwerth
aus den Werthen der Function $h$ zu beiden Seiten des Sprunges dar.
Bei der hier vorausgesetzten Bestimmungsweise der Function $f(x)$
besitzt diese dieselbe Eigenschaft, und man hat daher v\"{o}llig
allgemein
\[ f(y) = \frac{1}{2\pi i} \int\limits_{a - \infty i}^{a + \infty i}
      \frac{\log \zeta(s)}{s} y^s \,ds.\]

F\"{u}r $\log \zeta$ kann man nun den fr\"{u}her gefundenen Ausdruck
\[ \frac{s}{2} \log \pi - \log (s - 1)
   - \log \Pi \left( \frac{s}{2} \right)
   + {\textstyle\sum}^\alpha \log \left(
         1 + \frac{(s - \frac{1}{2})^2}{\alpha\alpha} \right)
   + \log \xi(0) \]
substituiren; die Integrale der einzelnen Glieder dieses Ausdrucks
w\"{u}rden aber dann ins Unendliche ausgedehnt nicht convergiren,
weshalb es zweckm\"{a}ssig ist, die Gleichung vorher durch partielle
Integration in
\[ f(x) = - \frac{1}{2\pi i} \frac{1}{\log x}
         \int\limits_{a - \infty i}^{a + \infty i} \, \,
         \frac{\displaystyle d \frac{\log \zeta(s)}{s}}{ds}
         x^s \,ds \]
umzuformen.

Da
\[ - \log \Pi \left( \frac{s}{2} \right)
   = \lim \left( \sum_{n=1}^{n = m}
         \log \left( 1 + \frac{s}{2n} \right)
         - \frac{s}{2} \log m \right),\]
f\"{u}r $m = \infty$, also
\[ - \frac{\displaystyle d \frac{1}{s}
            \log \Pi \left( \frac{s}{2} \right)}{ds}
   = \sum_1^\infty \frac{\displaystyle d \frac{1}{s}
            \log \left( 1 + \frac{s}{2n} \right)}{ds},\]
so erhalten dann s\"{a}mmtliche Glieder des Ausdruckes f\"{u}r $f(x)$
mit Ausnahme von
\[ \frac{1}{2\pi i} \frac{1}{\log x}
         \int\limits_{a - \infty i}^{a + \infty i}
         \frac{1}{ss} \log \xi(0) x^s \,ds
   = \log \xi(0) \]
die Form
\[ \pm \frac{1}{2\pi i} \frac{1}{\log x}
         \int\limits_{a - \infty i}^{a + \infty i} \, \,
         \frac{\displaystyle d \left( \frac{1}{s} \log \left(
               1 - \frac{s}{\beta} \right) \right)}{ds}
         x^s \,ds.\]

Nun ist aber
\[ \frac{\displaystyle d \left( \frac{1}{s} \log \left(
               1 - \frac{s}{\beta} \right) \right)}{d\beta}
   = \frac{1}{(\beta - s)\beta},\]
und, wenn der reelle Theil von $s$ gr\"{o}sser als der reelle Theil von
$\beta$ ist,
\[ - \frac{1}{2\pi i} \int\limits_{a - \infty i}^{a + \infty i}
         \frac{x^s \,ds}{(\beta - s)\beta}
   = \frac{x^\beta}{\beta} = \int\limits_\infty^x x^{\beta - 1} \,dx,\]
oder
\[ = \int\limits_0^x x^{\beta - 1} \,dx,\]
je nachdem der reelle Theil von $\beta$ negativ oder positiv ist.
Man hat daher
\[ \frac{1}{2\pi i} \frac{1}{\log x}
         \int\limits_{a - \infty i}^{a + \infty i} \, \,
         \frac{\displaystyle d \left( \frac{1}{s} \log \left(
               1 - \frac{s}{\beta} \right) \right)}{ds}
         x^s \,ds \]
\begin{eqnarray*}
   &=& - \frac{1}{2\pi i} \int\limits_{a - \infty i}^{a + \infty i}
         \frac{1}{s} \log \left( 1 - \frac{s}{\beta} \right)
         x^s \,ds \\
   &=& \int\limits_\infty^x \frac{x^{\beta - 1}}{\log x} \,dx
     + \mbox{ const.\ im ersten}
\end{eqnarray*}
und
\[
   =   \int\limits_0^x \frac{x^{\beta - 1}}{\log x} \,dx
     + \mbox{ const.\ im zweiten Falle.} \]

Im ersten Falle bestimmt sich die Integrationsconstante, wenn man den
reellen Theil von $\beta$ negativ unendlich werden l\"{a}sst; im
zweiten Falle erh\"{a}lt das Integral von $0$ bis $x$ um $2\pi i$
verschiedene Werthe, je nachdem die Integration durch complexe Werthe
mit positivem oder negativem Arcus geschieht, und wird, auf jenem Wege
genommen, unendlich klein, wenn der Coefficient von $i$ in dem Werthe
von $\beta$ positiv unendlich wird, auf letzterem aber, wenn dieser
Coefficient negativ unendlich wird.  Hieraus ergiebt sich, wie auf der
linken Seite
$\displaystyle \log \left( 1 - \frac{s}{\beta} \right)$
zu bestimmen ist, damit die Integrationsconstante wegf\"{a}llt.

Durch Einsetzung dieser Werthe in den Ausdruck von $f(x)$ erh\"{a}lt
man
\[ f(x) = {Li}(x) - {\textstyle\sum^\alpha} \left(
            {Li} \left( x^{\frac{1}{2} + \alpha i} \right)
            + {Li} \left( x^{\frac{1}{2} - \alpha i} \right)
            \right)
         + \int\limits_x^\infty \frac{1}{x^2 - 1} \frac{dx}{x \log x}
         + \log \xi(0),\]
wenn in $\sum^\alpha$ f\"{u}r $\alpha$ s\"{a}mmtliche positiven (oder
einen positiven reellen Theil enthaltenden) Wurzeln der Gleichung
$\xi(\alpha) = 0$, ihrer Gr\"{o}sse nach geordnet, gesetzt werden.  Es
l\"{a}sst sich, mit H\"{u}lfe einer genaueren Discussion der
Function~$\xi$, leicht zeigen, dass bei dieser Anordnung der Werth der
Reihe
\[ {\textstyle\sum} \left(
            {Li} \left( x^{\frac{1}{2} + \alpha i} \right)
            + {Li} \left( x^{\frac{1}{2} - \alpha i} \right)
            \right) \log x \]
mit dem Grenzwerth, gegen welchen
\[ \frac{1}{2\pi i} \int\limits_{a - bi}^{a + bi} \, \,
      \frac{\displaystyle d \frac{1}{s} {\textstyle\sum}
         \log \left( 1 + \frac{(s - \frac{1}{2})^2}{\alpha \alpha}
         \right)}{ds} x^s \,ds \]
bei unaufh\"{o}rlichem Wachsen der Gr\"{o}sse $b$ convergirt,
\"{u}bereinstimmt; durch ver\"{a}nderte Anordnung aber w\"{u}rde sie
jeden beliebigen reellen Werth erhalten k\"{o}nnen.

Aus $f(x)$ findet sich $F(x)$ mittelst der durch Umkehrung der
Relation
\[ f(x) = {\textstyle\sum} \frac{1}{n}
         F \left( x^{\frac{1}{n}} \right) \]
sich ergebenden Gleichung
\[ F(x) = {\textstyle\sum} (-1)^\mu \frac{1}{m}
         f \left( x^{\frac{1}{m}} \right),\]
worin f\"{u}r $m$ der Reihe nach die durch kein Quadrat ausser $1$
theilbaren Zahlen zu setzen sind und $\mu$ die Anzahl der Primfactoren
von $m$ bezeichnet.

Beschr\"{a}nkt man $\sum^\alpha$ auf eine endliche Zahl von Gliedern,
so giebt die Derivirte des Ausdrucks f\"{u}r $f(x)$ oder, bis auf
einen mit wachsendem $x$ sehr schnell abnehmenden Theil,
\[ \frac{1}{\log x} - 2 {\textstyle\sum^\alpha}
         \frac{ \cos (\alpha \log x) x^{-\frac{1}{2}}}{\log x} \]
einen angen\"{a}herten Ausdruck f\"{u}r die Dichtigkeit der Primzahlen
$+$ der halben Dichtigkeit der Primzahlquadrate $+$ $\frac{1}{3}$ von
der Dichtigkeit der Primzahlcuben u.~s.~w. von der Gr\"{o}sse $x$.

Die bekannte N\"{a}herungsformel $F(x) = {Li}(x)$ ist also nur bis auf
Gr\"{o}ssen von der Ordnung $x^{\frac{1}{2}}$ richtig und giebt einen
etwas zu grossen Werth; denn die nicht periodischen Glieder in dem
Ausdrucke von $F(x)$ sind, von Gr\"{o}ssen, die mit $x$ nicht in's
Unendliche wachsen, abgesehen:
\[ {Li}(x) - {\textstyle\frac{1}{2}}{Li}(x^{\frac{1}{2}})
           - {\textstyle\frac{1}{3}}{Li}(x^{\frac{1}{3}})
           - {\textstyle\frac{1}{5}}{Li}(x^{\frac{1}{5}})
           + {\textstyle\frac{1}{6}}{Li}(x^{\frac{1}{6}})
           - {\textstyle\frac{1}{7}}{Li}(x^{\frac{1}{7}})
           + \cdots \]

In der That hat sich bei der von \emph{Gauss} und
\emph{Goldschmidt} vorgenommenen und bis zu $x = $ drei Millionen
fortgesetzten Vergleichung von ${Li}(x)$ mit der Anzahl der Primzahlen
unter $x$ diese Anzahl schon vom ersten Hunderttausend an stets
kleiner als ${Li}(x)$ ergeben, und zwar w\"{a}chst die Differenz unter
manchen Schwankungen allm\"{a}hlich mit $x$.  Aber auch die von den
periodischen Gliedern abh\"{a}ngige stellenweise Verdichtung und
Verd\"{u}nnung der Primzahlen hat schon bei den Z\"{a}hlungen die
Aufmerksamkeit erregt, ohne dass jedoch hierin eine
Gesetzm\"{a}ssigkeit bemerkt worden w\"{a}re.  Bei einer etwaigen
neuen Z\"{a}hlung w\"{u}rde es interessant sein, den Einfluss der
einzelnen in dem Ausdrucke f\"{u}r die Dichtigkeit der Primzahlen
enthaltenen periodischen Glieder zu verfolgen.  Einen
regelm\"{a}ssigeren Gang als $F(x)$ w\"{u}rde die Function $f(x)$
zeigen, welche sich schon im ersten Hundert sehr deutlich als mit
${Li}(x) + \log \xi(0)$ im Mittel \"{u}bereinstimmend erkennen
l\"{a}sst.

\end{document}

